\begin{abstract}
    同策略蒸馏（OPD）将学生生成的轨迹与教师提供的稠密 token 级监督相结合。现有研究主要关注其算法行为，而对训练数据的作用则尚不明确。我们通过在单个查询上进行训练，考察了数据极小情况下的这一作用。one-shot OPD 在数百步内持续改进，并在任务领域和模型族之间恢复了大部分全数据 OPD 的收益。我们通过分析训练过程中访问的状态以及学生与教师对齐的速度来解释这一结果。我们测量了 \emph{状态覆盖}，即单个查询集的轨迹所达到的全数据 OPD 访问状态的比例。一个查询已达到 \(71.5\%\) ，其中大部分在前 100 步内完成。增加语义上不同的查询可同时提升覆盖度和验证准确率，直到 16 个查询达到 \(98.9\%\) ，并匹配全数据训练的表现。然而，无论 OPD 是基于单个查询还是整个数据集训练，对齐速度均以相似速率减慢，甚至固定的状态集合也需要数百步才能被吸收。因此，OPD 存在数据过载但算法饥饿的问题：其轨迹能迅速暴露广泛监督，而学生吸收这些监督的速度却越来越慢。状态覆盖的结果也适用于多教师 OPD，其中每个领域 16 个语义多样化的查询即可匹配全数据 MOPD。作为进一步的压力测试，内容稀疏的模板和跨领域的 WildChat 查询也接近真实查询基线。这表明任务内容与诱导出的状态覆盖可以分离。我们希望这些发现能引导未来工作关注 OPD 的步骤效率，并促使重新审视其在前沿后训练中取得近期成功的数据与机制基础。
\end{abstract}

